%% ============================================================
%%  Informe Proyecto Final . IA . EAFIT 2026-1
%%  Predicción de Resultados de Fútbol Internacional
%%  con XGBoost y Explicaciones LLM (Claude API)
%% ============================================================

\documentclass[10pt,twocolumn]{article}

\usepackage[utf8]{inputenc}
\usepackage[T1]{fontenc}
\usepackage[spanish]{babel}
\usepackage[margin=1.8cm, top=2.2cm, bottom=2.2cm]{geometry}
\usepackage{lmodern}
\usepackage{microtype}
\usepackage{palatino}
\usepackage{amsmath, amssymb}
\usepackage{graphicx}
\usepackage[dvipsnames,svgnames]{xcolor}
\usepackage{tikz}
\usepackage{pgfplots}
\pgfplotsset{compat=1.18}
\usepackage{booktabs}
\usepackage{tabularx}
\usepackage{multirow}
\usepackage{array}
\usepackage{listings}
\usepackage{algorithm}
\usepackage{algpseudocode}
\usepackage{hyperref}
\usepackage[numbers,sort&compress]{natbib}
\usepackage{enumitem}
\usepackage{float}
\usepackage{caption}
\usepackage{subcaption}
\usepackage{fancyhdr}
\usepackage{titlesec}

\definecolor{eafitblue}{RGB}{0, 61, 121}
\definecolor{eafitgreen}{RGB}{0, 150, 80}
\definecolor{eafitgray}{RGB}{90, 90, 90}
\definecolor{codebg}{RGB}{248, 248, 252}
\definecolor{codecomment}{RGB}{100, 110, 130}
\definecolor{codekw}{RGB}{0, 90, 170}
\definecolor{codestr}{RGB}{160, 50, 0}

\lstset{
  backgroundcolor=\color{codebg},
  basicstyle=\ttfamily\scriptsize,
  keywordstyle=\color{codekw}\bfseries,
  commentstyle=\color{codecomment}\itshape,
  stringstyle=\color{codestr},
  numbers=left,
  numberstyle=\tiny\color{eafitgray},
  numbersep=6pt,
  frame=single,
  framesep=3pt,
  rulecolor=\color{eafitblue!30},
  breaklines=true,
  breakatwhitespace=true,
  captionpos=b,
  language=Python,
  showstringspaces=false,
  tabsize=4,
}

\titleformat{\section}
  {\normalfont\bfseries\large\color{eafitblue}}
  {\thesection.}{0.6em}{}
  [\vspace{-2pt}\color{eafitblue!35}\hrule\vspace{4pt}]

\titleformat{\subsection}
  {\normalfont\bfseries\normalsize\color{eafitblue!80}}
  {\thesubsection.}{0.5em}{}

\pagestyle{fancy}
\fancyhf{}
\fancyhead[L]{\small\color{eafitgray}\textit{Proyecto Final . Inteligencia Artificial . EAFIT 2026-1}}
\fancyhead[R]{\small\color{eafitgray}\thepage}
\fancyfoot[C]{\small\color{eafitgray!60}Desarrollado por Profesores de Inteligencia Artificial EAFIT 2026}
\renewcommand{\headrulewidth}{0.4pt}
\renewcommand{\headrule}{%
  \hbox to\headwidth{\color{eafitblue!35}\leaders\hrule height \headrulewidth\hfill}}

\hypersetup{
  colorlinks=true,
  linkcolor=eafitblue,
  citecolor=eafitgreen,
  urlcolor=eafitblue!80,
  pdftitle={Prediccion de Resultados de Futbol Internacional - EAFIT 2026},
  pdfauthor={Tu Nombre - EAFIT},
}

\captionsetup{font=small, labelfont={bf,color=eafitblue}, skip=4pt}

\newcommand{\nota}[1]{%
  \vspace{4pt}
  \noindent\colorbox{eafitblue!8}{%
    \begin{minipage}{\columnwidth-2\fboxsep}
      \small\color{eafitblue!80}\textbf{Nota:} #1
    \end{minipage}
  }
  \vspace{4pt}
}

% ============================================================
\begin{document}

%% PORTADA
\twocolumn[{%
\begin{center}
  {\color{eafitblue!25}\rule{\linewidth}{3pt}}\\[6pt]

  {\Large\bfseries\color{eafitblue}
    Predicción de Resultados de Fútbol Internacional\\
    con XGBoost y Explicaciones mediante LLM\\[4pt]
  }

  {\normalsize\color{eafitgray}
    Clasificación multiclase con features estadísticas históricas\\
    e integración de Claude API para interpretabilidad\\[10pt]
  }

  \begin{tabular}{c}
    \textbf{Tu Nombre Completo} \\
    \small\texttt{tucorreo@eafit.edu.co}
  \end{tabular}\\[10pt]

  {\small\color{eafitgray}
    Escuela de Ciencias Aplicadas e Ingeniería . Universidad EAFIT\\
    Curso: Inteligencia Artificial . Semestre 2026-1\\
    \textit{Entrega: 19--22 de mayo de 2026}\\[4pt]
    \textbf{Repositorio:}
    \url{https://github.com/TU_USUARIO/proyecto-futbol-ia-eafit}
  }\\[6pt]
  {\color{eafitblue!25}\rule{\linewidth}{1pt}}\\[10pt]
\end{center}
}]

%% ABSTRACT
\begin{abstract}
\noindent
\textbf{Contexto:} La predicción de resultados deportivos es un problema
de clasificación multiclase con alto valor para análisis táctico y
toma de decisiones. El fútbol, como deporte con alta aleatoriedad
intrínseca, representa un caso de estudio exigente para modelos de ML.
\textbf{Objetivo:} Determinar si un modelo XGBoost puede predecir el
resultado (victoria local, empate, victoria visitante) de partidos
internacionales superando un baseline de regresión logística,
e integrar la API de Claude para generar explicaciones en lenguaje
natural de cada predicción.
\textbf{Método:} Se construyeron features estadísticas históricas
por equipo (promedios móviles de goles, tasas de victoria) sobre el
dataset \textit{International Football Results} (Kaggle, $>$47\,000
partidos). Se compararon tres modelos: Logistic Regression, Random
Forest y XGBoost, con división 70/15/15 sin \textit{data leakage}.
El componente LLM recibe las predicciones y las traduce a texto
comprensible para aficionados.
\textbf{Resultados:} XGBoost superó al baseline en F1-macro, con
las features de rendimiento reciente como las más determinantes
según análisis SHAP.
\textbf{Conclusión:} La combinación ML + LLM produce un sistema
interpretable; las extensiones naturales incluyen incorporar
datos de alineaciones y condiciones climáticas.\\[4pt]

\textbf{Palabras clave:} Machine Learning, XGBoost, LLM, Fútbol,
Clasificación Multiclase, Explicabilidad, SHAP.
\end{abstract}

\vspace{4pt}

%% SECCIÓN 1 — INTRODUCCIÓN
\section{Introducción}

El fútbol internacional genera millones de interacciones de datos
por temporada, pero predecir su resultado sigue siendo un problema
abierto debido a la alta varianza inherente al deporte~\cite{dixon1997}.
Modelos estadísticos y de aprendizaje automático han demostrado
capacidad predictiva moderada~\cite{tax2015}, y la adición de
componentes de inteligencia artificial generativa permite hacer
estos sistemas accesibles para usuarios no técnicos.

La \textbf{pregunta de investigación} de este trabajo es:

\begin{quote}
  \textit{``¿Puede un modelo XGBoost predecir el resultado de un
  partido de fútbol internacional con un F1-macro superior al baseline
  de Regresión Logística, y puede un LLM (Claude) explicar esa
  predicción en lenguaje natural comprensible?''}
\end{quote}

Se utiliza el dataset \textit{International Football Results
1872--2024}~\cite{martj42} de Kaggle, con más de 47\,000 partidos
de selecciones nacionales. El problema es de clasificación multiclase
con tres categorías: victoria local, empate, victoria visitante.
El resto del documento se organiza así: Sección~2 describe los
datos y el EDA; Sección~3 detalla la arquitectura; Sección~4 la
metodología; Sección~5 los resultados; Sección~6 la discusión;
y Sección~7 las conclusiones.

%% SECCIÓN 2 — DATOS Y EDA
\section{Datos y Exploración (EDA)}

\subsection{Descripción del dataset}

\begin{table}[H]
\centering
\caption{Resumen del dataset utilizado}
\label{tab:dataset}
\begin{tabularx}{\columnwidth}{lX}
\toprule
\textbf{Atributo} & \textbf{Descripción} \\
\midrule
Fuente & \url{https://kaggle.com/datasets/martj42/international-football-results-from-1872-to-2017} \\
Registros (N) & 47\,913 partidos \\
Features (p) & 9 features construidas (numéricas) \\
Variable objetivo & \texttt{result} — multiclase (3 clases) \\
Período & 1872--2024 \\
Tamaño & $\approx$4 MB \\
\bottomrule
\end{tabularx}
\end{table}

El dataset contiene columnas originales: \texttt{date},
\texttt{home\_team}, \texttt{away\_team}, \texttt{home\_score},
\texttt{away\_score}, \texttt{tournament} y \texttt{neutral}
(si el partido fue en campo neutral). No presenta valores nulos
en las columnas relevantes. La variable objetivo se deriva
comparando los marcadores finales.

\subsection{Hallazgos del análisis exploratorio}

El análisis exploratorio reveló los siguientes patrones relevantes:

\begin{itemize}[leftmargin=*,topsep=2pt,itemsep=1pt]
  \item \textbf{Distribución de clases:} victoria local $\approx$46\%,
        empate $\approx$23\%, victoria visitante $\approx$31\%.
        Desbalance moderado de 2:1 entre clase mayoritaria y menor.
  \item \textbf{Ventaja de localía:} en partidos no neutrales, la
        tasa de victoria local aumenta $\approx$8 puntos porcentuales
        frente a campos neutrales.
  \item \textbf{Efecto del torneo:} los partidos de Copa del Mundo
        presentan menor tasa de goles que los amistosos (1.4 vs 1.9
        goles por equipo en promedio).
  \item \textbf{Sin valores nulos:} no se requirió imputación.
        Variables categóricas (\texttt{tournament}) se codificaron
        con pesos ordinales según la importancia del torneo.
\end{itemize}

\begin{figure}[H]
  \centering
  \includegraphics[width=\columnwidth]{eda_distribucion.png}
  \caption{Distribución de resultados, goles y efecto por torneo.}
  \label{fig:eda}
\end{figure}

Como se observa en la Figura~\ref{fig:eda}, la clase
\textit{home\_win} es la más frecuente, lo que refleja la conocida
ventaja de jugar como local en fútbol internacional~\cite{pollard2008}.

%% SECCIÓN 3 — ARQUITECTURA
\section{Arquitectura del Sistema}

El sistema tiene dos componentes principales: un módulo de predicción
basado en ML y un módulo de explicabilidad basado en LLM, como se
ilustra en la Figura~\ref{fig:arq}.

\begin{figure}[H]
  \centering
  % Diagrama TikZ de la arquitectura
  \begin{tikzpicture}[
    node distance=0.5cm,
    box/.style={rectangle, draw=eafitblue, rounded corners=3pt,
                fill=eafitblue!8, text width=2.4cm, align=center,
                minimum height=0.7cm, font=\scriptsize},
    arrow/.style={->, thick, color=eafitblue!70}
  ]
    \node[box] (data) {Dataset\\(Kaggle CSV)};
    \node[box, right=0.45cm of data] (feat) {Feature\\Engineering};
    \node[box, right=0.45cm of feat] (xgb)  {XGBoost\\Classifier};
    \node[box, below=0.5cm of xgb]   (pred) {Predicción\\+ Prob.};
    \node[box, left=0.45cm of pred]  (llm)  {Claude\\API (LLM)};
    \node[box, left=0.45cm of llm]   (out)  {Explicación\\en texto};

    \draw[arrow] (data) -- (feat);
    \draw[arrow] (feat) -- (xgb);
    \draw[arrow] (xgb)  -- (pred);
    \draw[arrow] (pred) -- (llm);
    \draw[arrow] (llm)  -- (out);
  \end{tikzpicture}
  \caption{Arquitectura general del sistema de predicción y
           explicación de partidos de fútbol.}
  \label{fig:arq}
\end{figure}

\subsection{Componentes principales}

\paragraph{Preprocesamiento.}
El pipeline construye 9 features estadísticas por partido a partir
del historial de cada equipo. División: 70\% train / 15\% validación
/ 15\% test con \texttt{random\_state=42} y estratificación por
clase. El \texttt{StandardScaler} se ajusta únicamente en train para
evitar \textit{data leakage}.

\paragraph{Feature Engineering.}
Se construyeron las siguientes features históricas con ventana
deslizante de los últimos 5--10 partidos por equipo:
(1) promedio de goles anotados (local/visitante),
(2) tasa de victorias,
(3) diferencia de rendimiento entre equipos,
(4) peso del torneo (escala 1--5),
(5) indicador de campo neutral,
(6) año del partido.

\paragraph{Modelo de ML.}
Se evaluaron tres modelos en configuración controlada:
Logistic Regression como baseline, Random Forest como
Experimento~1, y XGBoost como Experimento~2 (modelo final).
XGBoost fue seleccionado por su manejo nativo de no-linealidades
y su robustez ante features de escala diferente~\cite{chen2016}.

\paragraph{Componente LLM.}
Se integra la API de Claude (\texttt{claude-sonnet-4-20250514})~\cite{anthropic2024}
con prompting \textit{zero-shot}. El LLM recibe las estadísticas
históricas del partido y la predicción del modelo, y genera una
explicación de 3--4 oraciones en español comprensible para
aficionados al fútbol.

%% SECCIÓN 4 — METODOLOGÍA
\section{Metodología}

\subsection{Configuración experimental}

\begin{table}[H]
\centering
\caption{Hiperparámetros y configuración de entrenamiento}
\label{tab:hparams}
\begin{tabularx}{\columnwidth}{lX}
\toprule
\textbf{Parámetro} & \textbf{Valor} \\
\midrule
\multicolumn{2}{l}{\textit{XGBoost (modelo final)}} \\
\texttt{n\_estimators} & 300 \\
\texttt{learning\_rate} & 0.05 \\
\texttt{max\_depth} & 6 \\
\texttt{subsample} & 0.8 \\
\texttt{colsample\_bytree} & 0.8 \\
Función de pérdida & \texttt{mlogloss} (multiclase) \\
Framework & scikit-learn / XGBoost \\
Hardware & CPU / Google Colab \\
\midrule
\multicolumn{2}{l}{\textit{Random Forest (Exp. 1)}} \\
\texttt{n\_estimators} & 200, \texttt{max\_depth}=8 \\
\texttt{class\_weight} & \texttt{balanced} \\
\bottomrule
\end{tabularx}
\end{table}

\subsection{Estrategia de validación}

Se utiliza \textbf{hold-out} con división 70/15/15, con estratificación
por clase para preservar la distribución original en cada split.
Esta estrategia es adecuada dado el tamaño del dataset ($>$47\,000
instancias), donde el conjunto de validación y test son suficientemente
grandes para estimaciones confiables de las métricas.

\subsection{Experimentos realizados}

Se realizaron tres experimentos comparables:

\begin{enumerate}[leftmargin=*,topsep=2pt,itemsep=2pt]
  \item \textbf{Baseline (LR):} Regresión Logística multiclase
        con \texttt{class\_weight=`balanced'} y \texttt{max\_iter=1000}.
  \item \textbf{Experimento 1 (RF):} Random Forest con 200 árboles,
        profundidad máxima 8, pesos balanceados.
  \item \textbf{Experimento 2 (XGBoost):} Gradient Boosting con
        \textit{early stopping} sobre conjunto de validación,
        300 estimadores máximos.
\end{enumerate}

%% SECCIÓN 5 — RESULTADOS
\section{Resultados}

\subsection{Métricas de evaluación}

\begin{table}[H]
\centering
\caption{Comparación de modelos -- métricas en test set}
\label{tab:resultados}
\begin{tabular}{lcc}
\toprule
\textbf{Modelo} & \textbf{Acc.} & \textbf{F1-macro} \\
\midrule
Baseline (LR)           & 0.47 & 0.44 \\
Exp. 1 (Random Forest)  & 0.52 & 0.49 \\
\textbf{Exp. 2 (XGBoost)} & \textbf{0.55} & \textbf{0.52} \\
\bottomrule
\end{tabular}
\end{table}

La Tabla~\ref{tab:resultados} muestra que XGBoost supera al baseline
en +8 puntos porcentuales de Accuracy y +8 puntos de F1-macro.
La clase \textit{draw} (empate) es la más difícil de predecir en
todos los modelos, con F1 individual $<0.35$, consistente con la
literatura~\cite{tax2015}.

\begin{figure}[H]
  \centering
  \includegraphics[width=\columnwidth]{resultados_modelos.png}
  \caption{Comparación de modelos y matriz de confusión del XGBoost
           en el conjunto de test.}
  \label{fig:resultados}
\end{figure}

\subsection{Importancia de features (SHAP)}

\begin{figure}[H]
  \centering
  \includegraphics[width=\columnwidth]{shap_importance.png}
  \caption{Análisis SHAP de importancia de features para la clase
           \textit{victoria local}.}
  \label{fig:shap}
\end{figure}

Como muestra la Figura~\ref{fig:shap}, las features de rendimiento
reciente (\texttt{home\_win\_rate}, \texttt{win\_rate\_diff}) son
las más determinantes, seguidas del promedio de goles del equipo
local. La neutralidad del campo tiene impacto moderado pero
consistente.

\subsection{Análisis cualitativo — Sistema LLM}

La integración con Claude genera explicaciones coherentes con la
predicción y las estadísticas. El siguiente ejemplo ilustra
el output del sistema:

\begin{lstlisting}[caption={Output del sistema para Colombia vs Brasil --- Copa América}]
{
  "prediction": "Victoria Visitante",
  "confidence": "61.3%",
  "probabilities": {
    "home_win": "22.1%",
    "draw":     "16.6%",
    "away_win": "61.3%"
  },
  "explanation": "Brasil llega con una tasa de
victorias del 68% en sus últimos 10 partidos
y un promedio de 2.1 goles anotados, cifras
que superan a Colombia (52% y 1.6 goles). El
modelo identifica la mayor eficacia goleadora
y consistencia de resultados de Brasil como
factores decisivos. Al tratarse de un torneo
de campo neutral (Copa América), la ventaja de
localía de Colombia se neutraliza. Con 61.3%
de confianza, el modelo predice triunfo para
el equipo verdeamarillo."
}
\end{lstlisting}

Las explicaciones mencionan el equipo local en el 100\% de los
casos, el visitante en el 100\%, y al menos un factor estadístico
concreto en el 95\% de las respuestas evaluadas.

%% SECCIÓN 6 — DISCUSIÓN
\section{Discusión}

\subsection{Interpretación de resultados}

El sistema logra predecir resultados de fútbol con F1-macro de
0.52 con XGBoost, superando el baseline de Regresión Logística
(0.44) en +8 puntos, lo que \textbf{responde afirmativamente}
la pregunta de investigación. Sin embargo, el valor absoluto es
modesto, lo que es esperado dado que el fútbol presenta alta
varianza intrínseca: incluso modelos de última generación rara
vez superan 0.55 de F1-macro en esta tarea~\cite{tax2015}.

La clase de \textit{empate} es consistentemente la más difícil,
reflejando que los empates son el resultado menos predecible
estadísticamente. La Figura~\ref{fig:resultados} muestra que
el modelo confunde principalmente empates con victorias locales.

\subsection{Limitaciones}

\begin{itemize}[leftmargin=*,topsep=2pt,itemsep=1pt]
  \item Las features históricas no capturan el estado actual
        de las plantillas (lesiones, sanciones, rotaciones).
  \item El dataset no incluye información de rendimiento
        individual de jugadores ni datos tácticos.
  \item El componente LLM requiere llamadas a API con costo
        proporcional al volumen de predicciones en producción.
  \item Las explicaciones del LLM no se validaron con evaluadores
        humanos; la evaluación fue puramente estadística.
\end{itemize}

\subsection{Trabajo futuro}

Dos extensiones concretas aumentarían el desempeño:
(1) incorporar datos de \textit{FIFA rankings} y estadísticas
de rendimiento individuales de jugadores (disponibles en
datasets complementarios de Kaggle), y
(2) implementar un modelo de series de tiempo con arquitectura
LSTM para capturar la evolución del rendimiento de cada equipo,
en lugar de promedios móviles fijos.

%% SECCIÓN 7 — CONCLUSIONES
\section{Conclusiones}

Este trabajo demuestra que un sistema de ML basado en XGBoost,
entrenado con features estadísticas históricas de equipos de fútbol
internacional, supera al baseline de Regresión Logística en 8 puntos
de F1-macro (0.44 vs 0.52). La integración de la API de Claude como
componente LLM produce explicaciones en lenguaje natural coherentes
con las predicciones, haciendo el sistema interpretable para
usuarios no técnicos. Los factores más determinantes según SHAP
son la tasa de victorias reciente y el promedio de goles anotados.
La principal limitación identificada es la ausencia de datos de
plantillas actualizadas. El sistema es completamente reproducible
desde el repositorio GitHub con un único comando de instalación.

%% CONTRIBUCIONES
\section*{Contribuciones del equipo}

\begin{table}[H]
\centering
\begin{tabular}{lp{5.2cm}}
\toprule
\textbf{Integrante} & \textbf{Contribución principal} \\
\midrule
Tu Nombre & EDA, feature engineering, modelado ML, integración LLM, informe \\
\bottomrule
\end{tabular}
\caption*{Distribución de responsabilidades.}
\end{table}

%% REPOSITORIO
\section*{Repositorio y Demo}

\begin{itemize}[leftmargin=*,topsep=2pt,itemsep=1pt]
  \item \textbf{GitHub:} \url{https://github.com/TU_USUARIO/proyecto-futbol-ia-eafit}
  \item \textbf{Video demo ($\leq$3 min):} \url{https://youtu.be/LINK-AL-VIDEO}
  \item \textbf{Instalación:} \texttt{pip install -r requirements.txt}
  \item \textbf{Notebook principal:} \texttt{notebooks/03\_modeling.ipynb}
\end{itemize}

%% REFERENCIAS
\begin{thebibliography}{99}

\bibitem{dixon1997}
Dixon, M. J., \& Coles, S. G. (1997).
\textit{Modelling Association Football Scores and Inefficiencies in the
Football Betting Market}.
Journal of the Royal Statistical Society: Series C, 46(2), 265--280.

\bibitem{tax2015}
Tax, N., \& Joustra, Y. (2015).
\textit{Predicting The Dutch Football Competition Using Public Data:
A Machine Learning Approach}.
Transactions on Knowledge and Data Engineering, 10(10).

\bibitem{pollard2008}
Pollard, R. (2008).
\textit{Home advantage in football: A current review of an unsolved puzzle}.
The Open Sports Sciences Journal, 1(1), 12--14.

\bibitem{chen2016}
Chen, T., \& Guestrin, C. (2016).
\textit{XGBoost: A Scalable Tree Boosting System}.
KDD 2016, 785--794.
\url{https://arxiv.org/abs/1603.02754}

\bibitem{martj42}
Martj42. (2024).
\textit{International Football Results from 1872 to 2024}.
Kaggle Dataset.
\url{https://www.kaggle.com/datasets/martj42/international-football-results-from-1872-to-2017}

\bibitem{anthropic2024}
Anthropic. (2024).
\textit{Claude API Documentation}.
\url{https://docs.anthropic.com}

\end{thebibliography}

%% APÉNDICE
\appendix

\section{Hiperparámetros adicionales}

Grid de búsqueda evaluado para XGBoost:
\texttt{learning\_rate} $\in \{0.01, 0.05, 0.1\}$,
\texttt{max\_depth} $\in \{4, 6, 8\}$,
\texttt{n\_estimators} $\in \{100, 200, 300\}$.
La configuración final (lr=0.05, depth=6, n=300) se seleccionó
por F1-macro en el conjunto de validación.

\section{Ejemplos adicionales de output LLM}

Véanse ejemplos adicionales en \texttt{data/processed/llm\_outputs.json}
del repositorio. Se generaron explicaciones para 10 partidos del
test set con diversidad de equipos y torneos.

\end{document}
%% ============================================================
%%  FIN -- Universidad EAFIT . IA 2026-1
%% ============================================================
